\documentclass[10pt,twocolumn,letterpaper]{article}

\usepackage{dicta}
\usepackage{times}
\usepackage{epsfig}
\usepackage{amsmath}
\usepackage{amssymb}
\usepackage{booktabs}
\usepackage{graphicx}
\usepackage{multirow}
\usepackage{xcolor}
\usepackage{placeins}
\usepackage{xfp}

% Keep hyperref after most packages, as requested by the DICTA template.
\usepackage[pagebackref=false,breaklinks=true,letterpaper=true,colorlinks,bookmarks=false]{hyperref}

% Toggle author visibility.
% Use \showauthorsfalse for the initial double-blind DICTA submission.
% Use \showauthorstrue for camera-ready or internal named drafts.
\newif\ifshowauthors
\showauthorsfalse
\ifshowauthors
\dictafinalcopy
\fi
% \dictafinalcopy
\def\dictaPaperID{****} % Replace with the DICTA paper ID once assigned.
\def\httilde{\mbox{\tt\raisebox{-.5ex}{\symbol{126}}}}
\newcommand{\mainfigwidth}{0.99\textwidth}
\newcommand{\widefigheight}{0.40\textheight}
\newcommand{\qualfigheight}{0.41\textheight}
\newcommand{\partialmetric}[1]{#1}
\newcommand{\bestpartialmetric}[1]{\underline{#1}}
\newcommand{\metricformat}[3]{%
  \ifnum\fpeval{abs((#1)-(#3)) < 0.00005}=1
    \ifnum\fpeval{#2}=1 \textbf{\underline{#1}}\else \underline{#1}\fi
  \else
    \ifnum\fpeval{#2}=1 \textbf{#1}\else #1\fi
  \fi}
\newcommand{\metriccd}[3]{\metricformat{#1}{(#1)<(#2)}{#3}}
\newcommand{\metricf}[3]{\metricformat{#1}{(#1)>(#2)}{#3}}
\newcommand{\metricbest}[2]{%
  \ifnum\fpeval{abs((#1)-(#2)) < 0.00005}=1 \underline{#1}\else #1\fi}
\newcommand{\pendingmetric}{--}
\newcommand{\resulttablesize}{\scriptsize}
% Keep result-table font sizes consistent; adjust only cell padding/row stretch.
\newcommand{\resulttablespacing}[1]{\setlength{\tabcolsep}{#1}\renewcommand{\arraystretch}{0.92}}
\makeatletter
\setlength{\@fptop}{0pt}
\setlength{\@dblfptop}{0pt}
\makeatother
\setlength{\textfloatsep}{5pt plus 1pt minus 1pt}
\setlength{\dbltextfloatsep}{5pt plus 1pt minus 1pt}
\setlength{\floatsep}{5pt plus 1pt minus 1pt}
\setlength{\dblfloatsep}{5pt plus 1pt minus 1pt}
\setlength{\intextsep}{5pt plus 1pt minus 1pt}
\setlength{\abovecaptionskip}{2pt}
\setlength{\belowcaptionskip}{0pt}
\renewcommand{\topfraction}{0.97}
\renewcommand{\dbltopfraction}{0.97}
\renewcommand{\textfraction}{0.03}
\renewcommand{\floatpagefraction}{0.82}
\renewcommand{\dblfloatpagefraction}{0.82}

\newcommand{\SK}[1]{\textcolor{blue}{[{\sc SK}] #1}}

% Pages are numbered/rulered in review mode and unnumbered in camera-ready.
\ifdictafinal\pagestyle{empty}\fi

\title{Curated Target-Part Preferences for Part-Aware Three-Dimensional (3D) Generation}

\author{
\begin{tabular}{c}
{\large Zishan Qin$^{1,2}$ \quad Rui Li$^{3}$ \quad Zeyu Zhang$^{4}$ \quad Lina Yao$^{1}$}\\
{\large Sevvandi Kandanaarachchi$^{2}$ \quad Dong Gong$^{1,*}$}\\[2pt]
$^{1}$University of New South Wales (UNSW Sydney), Australia\\
$^{2}$Commonwealth Scientific and Industrial Research Organisation (CSIRO)'s Data61, Australia\\
$^{3}$Nanyang Technological University, Singapore\\
$^{4}$University of California, Berkeley, United States of America\\[2pt]
{\tt\small \{zishan.qin,lina.yao,dong.gong\}@unsw.edu.au}\\
{\tt\small lirui.david@gmail.com, steve.zeyu.zhang@outlook.com, Sevvandi.Kandanaarachchi@csiro.au}\\
{\small $^*$Corresponding author}
\end{tabular}
}

\begin{document}
\maketitle
\begin{abstract}
% % Previous abstract: Part-level 3D generation aims to produce objects whose semantic components are not only recognizable, but also separated, complete, and editable.
% % Previous abstract: Existing 3D generation methods increasingly improve object-level fidelity, but they still lack effective supervision for local semantic parts, especially because semantically meaningful, human-aligned part annotations are scarce.
% % Previous abstract: Preference-based learning offers a new viewpoint: instead of asking annotators or automatic rules to provide a fully corrected 3D part, it can learn from comparisons between two pretrained-generator candidate target parts under the same input image.
% % Previous stronger framing: This paper studies whether such preference supervision can provide a useful learning signal for improving weak semantic parts in part-aware 3D generation.
% % Previous abstract: This paper studies whether such preference supervision can provide a useful learning signal for part-level 3D generation when fully corrected part-level targets are costly.
% % Previous abstract: Specifically, we curate an inspectable candidate set generated by the pretrained PartCrafter model on PartObjaverse-Tiny source objects, containing 300 same-input target-part preference examples.
% % Previous abstract: We define preference construction criteria for metric based, ground-truth-guided selection, synthetic negative construction, and human inspection, then train with a reference-anchored pairwise objective plus a denoising regularizer.
% % Previous abstract: The experiments examine five factors: pair-construction strategy, positive--negative quality, effective pair-set size, cross-dataset transfer, and robustness to noisy or synthetic preferences.
% % Previous abstract: Current results show that the ground-truth-guided preference setting improves the closest held-out Objaverse multi-part Part F-score from 0.55 to 0.57 and reduces Part CD from 0.29 to 0.27 over the pretrained baseline.
% % Previous abstract: These findings support a conservative conclusion: curated part-level preferences can provide a positive part-level supervision signal, but the benefit depends on preference quality, negative construction, data coverage, objective balance, and dataset shift.
% Although recent 3D generation methods have improved object-level fidelity, plausible global shape does not guarantee correct local semantic structure. This gap motivates part-aware generation, where semantic components should be separated, complete, and individually editable. 
% However, direct part-level supervision is costly because it requires identifying semantic components that may be fused in generated shapes and aligning part decompositions with human interpretations.
% % However, collecting direct part-level supervision remains difficult because it requires identifying meaningful components, verifying local geometry, and aligning part decompositions with human semantic understanding.
% We therefore study whether pairwise preferences can reduce this supervision burden by learning from comparisons between generated alternatives rather than requiring fully corrected target-part meshes for every example. Rather than relying on dense ground-truth part annotations during training, our approach learns from pairwise comparisons between two candidate target parts generated by a pretrained part-aware model under the same input image and part condition. Using a small curated pool of approximately 800 part-level pairs, each preference example contrasts a preferred part mesh with a weaker alternative.
% % Previous abstract wording: We construct preferences through metric/ground-truth-assisted automatic pair selection and human inspection, and study additional controls including synthetic negative construction and controlled pair-order perturbations.
% % Previous wording: We systematically evaluate the effects of pair-construction strategy, effective pair-set size, and cross-dataset transfer.
% We construct preferences through metric-assisted automatic selection and human inspection, and study additional controls including synthetic negative construction, pair-order perturbation, and uncurated noisy-sample injection. The model is optimized with a reference-anchored pairwise objective combined with a denoising regularizer. We systematically evaluate the effects of pair source, effective pair-set size, label reliability, negative construction, and cross-dataset transfer. In the strongest fully supported metric-assisted Objaverse variant, the preference-supervised model improves Part F-score (FS) from 0.55 to 0.57 and reduces Part Chamfer Distance (CD) from 0.29 to 0.25 compared with the pretrained baseline. 
% % Control results show that randomly ordered noisy samples do not replace curated preferences. 
% These findings support 
% % a conservative 
% the conclusion that curated part-level preferences can serve as a useful supervision signal for part-aware generation, while the benefit depends on preference quality, data scale, and objective balance ratio.

Although recent three-dimensional (3D) generation methods have improved object-level fidelity, plausible global shape does not guarantee correct local semantic structure. This gap motivates part-aware generation, where semantic components should be separated, complete, and individually editable. However, direct part-level supervision is costly because generated parts may be fused or ambiguous, requiring component identification and alignment with human semantic interpretations.
We therefore study whether pairwise preferences can reduce the cost of training supervision by learning from a compact set of comparisons between generated target-part samples rather than requiring fully corrected target-part meshes for every example.
% We therefore study whether pairwise preferences can reduce this supervision burden by learning from comparisons between generated target-part alternatives rather than requiring fully corrected target-part meshes for every example.
% Previous visible wording: Given the same input image and target part, we compare two candidates produced by pretrained part-aware generation models and train the denoising transformer with a reference-anchored pairwise objective and a denoising regularizer. We use a 300-sample metric-assisted preference set, evaluate human inspection as a smaller annotation-source pilot, and study pair-set size, label reliability, synthetic negatives, noisy-sample injection, and cross-dataset transfer. In the strongest metric-assisted Objaverse setting, curated preferences improve Part F-score (FS) from 0.55 to 0.57 and reduce Part Chamfer Distance (CD) from 0.29 to 0.25 over the pretrained baseline. These results show that curated part-level preferences can improve part-aware generation, while the gain depends on preference quality, data coverage, and objective weighting between local preference adaptation and object-level preservation.
Given the same input image and target part, we form preference pairs of generated target-part samples and train the denoising transformer with a DPO-style pairwise objective relative to a frozen reference model plus a denoising regularizer. We use a 300-sample metric-assisted preference set, evaluate human inspection through two small annotation-source pilots, and study pair-set size, label reliability, synthetic negatives, noisy-sample injection, and cross-dataset transfer. In the strongest metric-assisted Objaverse setting, curated preferences improve Part F-score (FS) from 0.55 to 0.57 and reduce Part Chamfer Distance (CD) from 0.29 to 0.25 over the pretrained baseline. These results show that curated part-level preferences can improve part-aware generation, while the gain depends on preference quality, data coverage, and objective weighting between local preference adaptation and object-level preservation.
\end{abstract}

\section{Introduction}
Recent image-to-3D and text-to-3D methods can synthesize plausible 3D objects from image or text conditions \cite{poole2022dreamfusion,liu2023zero1to3,hong2023lrm}. However, global plausibility does not guarantee part-level semantic correctness: local components must be identifiable and separated for editing, recomposition, or articulation \cite{mo2019partnet,wang2019shape2motion,li2026partrag,liu2024singapo,liu2026pact}. A generated object may preserve a recognizable silhouette while subtle parts are missing, fused, or weakly separated.

% Previous introduction: This work focuses on the supervision gap behind this failure mode. Part-level reconstruction data are more expensive to collect and verify than object-level examples: one must identify semantic components, ensure local geometric validity, and avoid ambiguous labels for tiny or weakly visible parts. As a result, a part-aware generator can receive enough signal to produce a reasonable object while still receiving limited direct supervision for local semantic components.
This work focuses on the supervision gap behind these failures. Direct part-level supervision is costly because generated components may be fused or ambiguous, requiring component identification and alignment with human semantic interpretations. A generator may therefore learn recognizable global shape while receiving limited evidence about which local parts should be complete, separated, and editable.
% Previous stronger framing: We therefore ask whether preference-based supervision can provide a lighter way to improve weak semantic parts.
We therefore study whether pairwise preferences can reduce the cost of obtaining training supervision by learning from comparisons between generated target-part samples rather than requiring fully corrected target-part meshes for every example.

% Previous introduction: The study is framed around the part-level reconstruction task rather than a method specialized to one generator. We use PartCrafter \cite{lin2025partcrafter} as the experimental backbone because it already produces multiple part meshes from a single input image, but the core question is broader: can same-input, same-target-part preferences provide a useful learning signal when full part-level targets are costly? A preferred branch and a rejected branch are comparable because they share the same input image and target semantic part. The preference label can be obtained from metric/ground-truth-assisted selection, synthetic negative construction, or human inspection.
% Previous wording: The preference label can be obtained from metric-assisted generated-pair selection, synthetic negative construction, controlled perturbation, or human inspection.
% Previous visible wording: We use PartCrafter \cite{lin2025partcrafter} as the experimental backbone because it generates multiple part meshes from a single input image. The study is not intended as a new generator architecture, but as an evaluation of whether same-input, same-target-part preferences provide useful local supervision. A preferred branch and a rejected branch are comparable because they share the same input image and target semantic part. We distinguish pair-construction sources from reliability controls: sources define how preferred and rejected meshes are selected or constructed, while controls test pair-set size, label perturbation, synthetic negatives, and objective strength.
We use PartCrafter \cite{lin2025partcrafter} as the experimental backbone because it already produces multiple part-level samples from a single input image. Rather than proposing a new generator architecture, we analyze same-input, same-target-part preference supervision. Comparing preferred and rejected target-part samples under fixed input and part conditions keeps the supervision compact, auditable, and focused on local quality.
% Previous stronger framing: This makes the supervision compact and auditable while directly targeting local part failures.
% Previous visible wording: This makes the supervision compact, auditable, and explicitly local to a target semantic part.

% Previous introduction: We present a controlled study using preference candidates generated from PartObjaverse-Tiny source objects \cite{yang2024sampart3d}. The primary metric-assisted setting contains 300 same-input target-part preference examples. We evaluate how preference learning behaves under different effective pair-set sizes, preference strengths, noisy labels, negative-sample construction strategies, and dataset shifts. The goal is not to claim a complete scalable annotation pipeline, nor to show that pairwise preference learning universally dominates supervised fine-tuning. Rather, the goal is to establish whether manually supervised part-level preferences can provide a positive signal, and to identify which data and objective factors control that signal.
% Previous wording: We evaluate how preference learning behaves under different effective pair-set sizes, pair-construction rules, noisy or synthetic preferences, preference strengths, and dataset shifts.
% Previous wording: We evaluate pair source, pair-set size, label reliability, synthetic-negative strength, preference strength, part-count sensitivity, and cross-dataset transfer.
% Previous visible wording: We conduct a controlled study using preference candidates generated from PartObjaverse-Tiny \cite{yang2024sampart3d}. The primary metric-assisted setting contains 300 same-input target-part preference examples; human inspection is evaluated as a smaller annotation-source pilot. We vary pair construction, pair-set size, label perturbation, synthetic rejected parts, uncurated noisy samples, objective weight, part count, and transfer dataset. The goal is not to claim a complete scalable annotation pipeline, but to test whether curated part-level preferences provide useful local supervision and which data or objective factors control the signal.
We conduct a controlled study using target-part samples generated from PartObjaverse-Tiny \cite{yang2024sampart3d}. The primary metric-assisted setting contains 300 same-input target-part preference examples; human inspection is evaluated through two small annotation-source pilots. We vary pair construction, pair-set size, label perturbation, synthetic rejected parts, uncurated noisy samples, objective weight, part count, and transfer dataset. The goal is not to claim a complete scalable annotation pipeline, but to test whether curated part-level preferences provide useful local supervision and which data or objective factors control the signal.

Our contributions are:
\begin{itemize}
    % Previous stronger framing: We formulate weak semantic parts as a data-supervision problem for part-level 3D reconstruction and generation.
    \item We formulate part-aware 3D generation as a supervision problem and study same-input, same-target-part preferences as a lightweight source of training signal over generated parts.
    % Previous contribution wording: We build an inspectable curated preference-candidate pool and define construction criteria for metric/ground-truth-assisted, synthetic-negative, perturbed, and human-inspected preferences.
    % Previous wording: We build an inspectable curated preference-candidate pool and define construction criteria for metric-assisted generated-pair, synthetic-negative, pair-perturbed, and human-inspected preferences.
    \item We build an inspectable preference-pair set from pretrained PartCrafter samples and define metric-assisted, human-inspected, and synthetic-negative pair-construction sources.
    % Previous stronger framing: We evaluate a reference-anchored pairwise preference objective as a lightweight supervision signal for weak semantic parts.
    % Previous wording: We evaluate how curated preferences behave as pair-set size, label reliability, negative construction, preference-objective weight, part count, and evaluation dataset change.
    \item We analyze how pair-set size, label noise, negative-pair construction, objective weight, target-part count, and evaluation dataset affect preference-supervised training.
\end{itemize}

% Previous visible wording: Figures~\ref{fig:preference_data_process}, \ref{fig:metric_pair_example}, and~\ref{fig:dit_preference_training} summarize the supervision loop: candidate-pool construction, same-input target-part pair selection, and reference-anchored DiT fine-tuning.
Figures~\ref{fig:preference_data_process} and~\ref{fig:metric_pair_example} summarize the main-paper supervision loop: sample generation and same-input target-part preference-pair selection. The supplementary material includes the DPO-style DiT fine-tuning schematic.

% Previous single-panel method overview. Replaced by
% Figs.~\ref{fig:metric_pair_example} and~\ref{fig:dit_preference_training},
% which separate preference-pair construction from DiT preference fine-tuning.
% \begin{figure*}[t]
% \centering
% \includegraphics[width=\mainfigwidth,height=\widefigheight,keepaspectratio]{figures/layout_consistent_20260630/method_overview_scientific_palette_cropped.png}
% \caption{Reference-anchored part-level preference learning for
% PartCrafter \cite{lin2025partcrafter}. Panel A isolates local part quality by keeping the input image and target semantic part fixed while comparing candidate branches generated from different stochastic seeds of the pretrained PartCrafter model. Target-part quality measures and curation checks select a preferred target-part mesh $y_p^+$ and a rejected target-part mesh $y_p^-$, so the preference label reflects a measured local part gap rather than a whole-object judgment. Panel B applies the resulting pair record to the released PartCrafter backbone: the trainable denoising model is encouraged to score $y_p^+$ above $y_p^-$ relative to a frozen pretrained reference, while the denoising regularizer on the preferred target part limits drift from the pretrained generator.}
% \label{fi为什么g:method}
% \end{figure*}

\section{Related Work}
To position the study, we review only the work needed to separate our question from prior advances in 3D generation: object-level generation, part-aware generation, editable assets, and preference alignment. A fuller related-work discussion is included in the supplementary material.

\noindent\textbf{Image-to-3D and text-to-3D generation.}
Recent 3D generation methods synthesize object-level geometry from text, images, or sparse views. Score-distillation methods use pretrained two-dimensional (2D) diffusion guidance \cite{poole2022dreamfusion,lin2023magic3d}, while sparse-view, feed-forward, and large-scale pipelines improve single-image or few-view reconstruction \cite{liu2023zero1to3,hong2023lrm,li2025triposg}. These systems provide strong object-level baselines, but their objectives do not guarantee separated, complete, or editable semantic parts.
% Previous stronger framing: Our work therefore focuses on the complementary question of how to improve weak semantic parts once a part-aware generator is already available.
Our work therefore focuses on the complementary question of how to obtain part-level supervision once a part-aware generator is already available.

\noindent\textbf{Part-aware 3D generation.}
Part-aware generation includes decompose-then-reconstruct pipelines that recover parts from segmentation, prompts, boxes, or multi-view cues \cite{chen2025partgen}; native part-generation methods that model parts as latent or token-level objects \cite{lin2025partcrafter,dong2025from,chen2025autopartgen}; and systems emphasizing structural layout, variable part count, or retrieval-augmented editing \cite{yang2025omnipart,tang2025efficient,li2026partrag}. These works primarily improve architecture, representation, or conditioning. We instead keep the pretrained backbone fixed and study curated same-input preferences as an additional supervision signal.

\noindent\textbf{Editable and articulated assets.}
A related line studies whether generated parts are functional, editable, or usable in downstream asset pipelines. Articulated-generation methods infer part graphs, joints, axes, or motion ranges \cite{liu2024singapo,liu2026pact}. These directions motivate our focus on local quality: missing, fused, or weak parts leave little reliable structure for later editing or articulation.

\noindent\textbf{Preference and reward alignment.}
Preference learning improves generative models when direct likelihood supervision is limited. Direct Preference Optimization (DPO) learns from pairwise preferences relative to a reference model \cite{rafailov2023direct}, and diffusion variants adapt the idea to denoising models \cite{wallace2024diffusiondpo}. Our implementation is a DPO-style pairwise objective with a frozen reference model, not an online reward-model or policy-gradient loop. Prior work on noisy preferences motivates our pair-order perturbation study \cite{raychowdhury2024robustdpo}. In 3D generation, preference feedback, human-preference alignment, and mesh refinement have been explored mostly at whole-object, rendered-view, or generic mesh-quality levels \cite{li2026circulardpo,zhou2025dreamdpoaligningtextto3dgeneration,liu2025meshrftenhancingmeshgeneration}.
% Previous stronger framing: Our study narrows the preference signal to same-input, same-target-part comparisons so that the update explicitly targets weak semantic parts rather than global asset quality alone.
Our study narrows the preference signal to same-input, same-target-part comparisons so that the supervision describes local target-part quality rather than global asset quality alone.

\section{Preliminary: Part-Level Generation Backbone}
To make the preference notation precise, we first fix the part-level generation backbone and identify which components are kept frozen. This section defines the input image $x$, target part $p$, sampled target-part meshes, and the frozen reference model used later in the method.

We use PartCrafter \cite{lin2025partcrafter} as the part-level generation backbone. Given a single input image $x$, it produces multiple part meshes rather than a single monolithic surface. The image encoder provides visual conditioning features, a PartCrafter Diffusion Transformer (DiT) denoiser predicts part-latent updates, and a TripoSG variational autoencoder (VAE) decoder \cite{li2025triposg} maps part latents into implicit geometry fields for mesh extraction. One stochastic generation run yields a full set of generated object parts; preference construction later selects the target semantic part for comparison.

We denote a target part by $p$, its ground-truth (GT) target mesh by $\mathrm{GT}_p$ when available, and a generated target-part mesh by $y_p$. The reference mesh is used only for pair construction or validation, not as a generator input. Let $\pi_{\mathrm{ref}}$ be the frozen pretrained generator. For each fixed $(x,p)$, we sample outputs from $\pi_{\mathrm{ref}}$ and extract $\mathcal{C}_{x,p}^{\mathrm{ref}}=\{y_{p,k}\}_{k=1}^{K}$; preference supervision is assigned to these extracted target-part samples.

The contribution is not a new generator architecture: the image encoder, VAE decoder, sampling scheduler, and backbone remain fixed while the denoising transformer is fine-tuned. For a fixed $(x,p)$ pair, $y_p^+$ and $y_p^-$ denote the preferred and rejected target-part samples from $\mathcal{C}_{x,p}^{\mathrm{ref}}$. This same-input construction makes the label a local part comparison rather than a broad whole-object judgment.
% Previous compact preliminary wording: Given a single input image $x$, the generator produces a set of part-conditioned 3D outputs rather than a single monolithic surface. The contribution of this paper is not a new part-aware generator architecture; instead, we study how to construct and use preference supervision over candidate target-part meshes generated by the pretrained backbone.
% \SK{I'd like to see an example with $x$  $p$, $y^+$ and $y^{-}$ are clearly shown.}

\section{Method}
To test whether compact part-level preferences can supervise weak local geometry, the method has three parts. Sec.~\ref{sec:pair_construction} defines same-input target-part pair construction, Sec.~\ref{sec:preference_objective} converts these pairs into a DPO-style objective relative to a frozen reference model, and Sec.~\ref{sec:experiment_protocol} describes the sampled-output set, pair-count controls, and evaluation protocol.

\subsection{Problem Setup}
Given input image $x$ and target part $p$, the goal is to improve local geometry while preserving surrounding object structure. A preference example is $(x,p,\mathrm{GT}_p,y_p^+,y_p^-)$ when reference geometry is available, or $(x,p,y_p^+,y_p^-)$ otherwise. Both samples come from $\mathcal{C}_{x,p}^{\mathrm{ref}}$ for the same input and target part. The label states only that $y_p^+$ is preferred to $y_p^-$ for that target part under the chosen annotation rule.

\subsection{Preference Pair Construction}
\label{sec:pair_construction}
% Previous wording: We organize preference construction into four variants. This organization separates the paper's central question from a single annotation source and makes the role of annotation quality explicit.
% Previous visible wording: We describe preference construction as three pair sources plus controls that test supervision reliability and coverage. Pair-set size, label flips, and crop ratio are therefore controls on the signal, not separate preference-learning methods.
We construct preferences to isolate local target-part quality while holding the input image and target part fixed. The primary source is GT/metric-assisted generated-pair selection; GT geometry or part-level metrics are used only to choose the preferred and rejected samples before training, not to define an additional supervising model. Human-inspected pairs and synthetic negatives are used as annotation-source and negative-construction controls. Pair count, label flips, noisy-sample injection, and crop ratio are treated as experimental controls on reliability and coverage rather than as separate methods.

\noindent\textbf{Sample generation and local comparison principle.}
For each retained preference example, we fix $x$ and $p$, then sample multiple outputs from the frozen pretrained model. Although each output may contain a full object, the preference label is assigned only to the extracted target-part mesh. The preferred sample is more complete, less artifact-prone, or better separated; the rejected sample keeps comparable object context but has weaker target-part geometry, such as missing structure, fusion, deformation, or artifacts. Fixing $x$ and $p$ makes the label a local reconstruction comparison rather than a category-level or whole-object preference.

\begin{figure*}[t]
\centering
\includegraphics[width=0.92\textwidth,keepaspectratio]{figures/figure1_pair_selection_chair_20260707/data_process_schematic/figure1_data_process_triposg_style_draft.png}
\caption{Preference-data construction pipeline. Source objects and target parts are rendered as input conditions, sampled with the frozen pretrained model, filtered, ranked by the part-level score $S_{\mathrm{auto},\mathrm{part}}$, and converted into same-input preference records. Human inspection and reliability controls are treated as separate variants.}
\label{fig:preference_data_process}
\end{figure*}

% Previous heading: Metric/ground-truth-assisted preferences.
\noindent\textbf{Metric-assisted generated-pair preferences.}
% Previous wording: When reference part geometry or reliable part-level metrics are available, candidate target-part meshes are ranked by target-part quality and filtered by curation checks. We use target-part Chamfer Distance (CD) and F-score (FS) as geometric ranking signals: CD measures average surface discrepancy, while FS measures the fraction of predicted and reference points matched under a distance threshold. The preferred mesh $y_p^+$ is selected from stronger target-part candidates, and the rejected mesh $y_p^-$ is selected from candidates with visibly weaker local geometry, missing parts, or poorer separation. This is the main source for the 796-effective-pair preference pool. The phrase \emph{metric-assisted} is used deliberately: metrics assist pair ordering and filtering, but the training example remains a generated preferred and rejected pair from the pretrained model rather than a dense supervised target.
% Previous visible detailed scoring block moved to the supplementary material:
% When reference geometry or reliable part-level metrics are available, an automatic target-part score ranks candidate branches before curation. We write it as $S_{\mathrm{auto},\mathrm{part}}$ because it selects preferred target-part meshes, not whole-object winners. Object-level terms act only as consistency checks against globally invalid branches. The score combines object consistency, part geometry, and lightweight structural checks:
% \begin{equation}
% \small
% \begin{aligned}
% S_{\mathrm{auto},\mathrm{part}}(y) =
% &w_{\mathrm{obj}}S_{\mathrm{obj}}(y)
% +w_{\mathrm{part}}S_{\mathrm{part}}(y)\\
% &+w_{\mathrm{aux}}S_{\mathrm{aux}}(y).
% \end{aligned}
% \label{eq:auto_score}
% \end{equation}
% Here $S_{\mathrm{obj}}$ aggregates object-level FS and CD checks, $S_{\mathrm{part}}$ aggregates part-level FS and CD terms, and $S_{\mathrm{aux}}$ includes normal consistency and component quality. Normal consistency measures surface-orientation agreement with reference geometry; component quality penalizes fragmented predicted parts using connected-component statistics. When multiple valid predicted parts are scored, part-level terms are averaged before applying the part-level weight:
% \begin{equation}
% \small
% \begin{aligned}
% S_{\mathrm{obj}}(y)&=
% \alpha_{\mathrm{obj}}^{\mathrm{FS}}\mathrm{FS}_{\mathrm{obj}}
% +\alpha_{\mathrm{obj}}^{\mathrm{CD}}(1-\mathrm{CD}_{\mathrm{obj}}),\\
% S_{\mathrm{part}}(y)&=
% \frac{1}{|\mathcal{P}_{\mathrm{valid}}|}
% \sum_{q\in\mathcal{P}_{\mathrm{valid}}}
% \Bigl[
% \alpha_{\mathrm{part}}^{\mathrm{FS}}\mathrm{FS}_{q}\\
% &\quad+\alpha_{\mathrm{part}}^{\mathrm{CD}}
% (1-\min(\mathrm{CD}_{q},1))
% \Bigr],\\
% S_{\mathrm{aux}}(y)&=
% \alpha_{\mathrm{NC}}\mathrm{NC}
% +\alpha_{\mathrm{CQ}}\mathrm{CQ}.
% \end{aligned}
% \end{equation}
% Coefficients are fixed across metric-assisted ranking runs and used only for candidate ordering, with object, part, and auxiliary terms normalized within groups before aggregation. Boundary-edge ratio and part intersection-over-union were retained only as calibration diagnostics and excluded from Eq.~\ref{eq:auto_score} because their ranking power was unstable. Higher $S_{\mathrm{auto},\mathrm{part}}$ indicates a better target-part candidate; $y_p^+$ is selected from high-scoring valid candidates and $y_p^-$ from lower-scoring candidates that pass non-empty and same-input checks. The training objective consumes only the resulting preferred--rejected ordering. Figure~\ref{fig:metric_pair_example} shows one candidate-ranking example.
When reference geometry or reliable part-level metrics are available, sampled outputs are ranked by a target-part score $S_{\mathrm{auto},\mathrm{part}}$ before curation. The score combines object-level consistency checks, target-part CD/FS terms, and lightweight normal/component diagnostics; its coefficients are used only for pair selection, not for training. We select $y_p^+$ from high-scoring valid target-part samples and $y_p^-$ from lower-scoring samples that pass non-empty and same-input checks. Full scoring details are provided in the supplementary material.

\begin{figure}[t]
\centering
\includegraphics[width=0.88\linewidth]{figures/metric_assisted_pair_example.png}
\caption{Preference-pair construction. For a fixed input and target part, sampled outputs are ranked by $S_{\mathrm{auto},\mathrm{part}}$; high- and low-scoring valid samples become $y_p^+$ and $y_p^-$. Intermediate samples make the automatic selection auditable, while training uses only the ordering.}
\label{fig:metric_pair_example}
\end{figure}

\noindent\textbf{Human-inspected preferences.}
% Previous draft wording: In the planned annotation-quality study, a human-inspected subset of approximately 100 examples will be compared against the metric/ground-truth-assisted rule to quantify agreement and downstream model behavior.
% Previous visible wording: Human inspection can directly choose preferred and rejected target parts or validate automatically mined pairs. It is useful because small components and ambiguous boundaries are not always captured by a single metric. We report the current human-inspected subset as an annotation-quality pilot.
Human inspection can directly choose preferred and rejected target-part samples or validate automatically mined pairs. We report two small annotation-source pilots: a 41-pair subset and seven ShapeNet pairs from five source objects, visually retained after metric and multiview screening. Because their source data and training settings differ, these pilots test annotation quality rather than pair-count scaling.

\noindent\textbf{Synthetic negative preferences.}
To reduce dependence on ground-truth part geometry, a preferred target-part sample can be converted into a synthetic negative by removing, masking, or corrupting the target part while keeping object context comparable. This mirrors occlusion-style image augmentations such as random erasing, Cutout, and Hide-and-Seek, but operates on 3D target-part geometry. Our \emph{+crop} variant removes a contiguous surface region along the longest spatial axis and fills removed slots by repeated kept points so the VAE input shape remains fixed. Removal ratios such as 30\% or 50\% define negative-construction strength.
% Previous stronger framing: This setting tests whether preference learning can benefit from synthetic weak-part negatives even when an exact GT-ranked rejected mesh is unavailable.
This setting tests whether preference learning can benefit from synthetic target-part negatives even when an exact ground-truth-ranked rejected mesh is unavailable.

\noindent\textbf{Experimental variables.}
After constructing pairs, we vary the supervision signal. Pair-set size tests coverage; pair-order flips test label reliability; noisy-sample injection adds uncurated seed-pair samples; synthetic-negative ratio tests handcrafted rejected targets; and preference strength with regularizer weights tests the balance between local adaptation and pretrained-generator preservation.

% Optional future extension: a purely heuristic negative-mining rule that does not use GT geometry can be added as an additional construction baseline if time permits.

\subsection{DPO-Style Pairwise Objective}
\label{sec:preference_objective}
% \SK{Was DPO previously defined?}
Direct Preference Optimization (DPO) compares a trainable policy against a frozen reference policy using pairwise preferences \cite{rafailov2023direct}. We use this form as the conceptual starting point because it asks whether a preferred sample should receive a higher implicit reward than a rejected sample relative to the reference policy for the same condition. For a conventional conditional generative policy and preference pair $(x,p,y_p^+,y_p^-)$, with $y_p^+,y_p^- \in \mathcal{C}_{x,p}^{\mathrm{ref}}$, the pairwise loss is
{\small
\begin{equation}
\begin{aligned}
\mathcal{L}_{\mathrm{DPO}}
= - \log \sigma \Bigg(
\beta \Big[
&\log \frac{\pi_\theta(y_p^+|x,p)}{\pi_{\mathrm{ref}}(y_p^+|x,p)} \\
&- \log \frac{\pi_\theta(y_p^-|x,p)}{\pi_{\mathrm{ref}}(y_p^-|x,p)}
\Big]
\Bigg),
\end{aligned}
\label{eq:dpo}
\end{equation}
}
where $\pi_\theta$ is the fine-tuned policy, $\pi_{\mathrm{ref}}$ is the frozen pretrained PartCrafter reference policy that also supplies the offline samples before curation, and $\beta$ controls the preference strength.

Because PartCrafter is a latent diffusion model, Eq.~\ref{eq:dpo} is not an exact tractable mesh likelihood. Following diffusion preference alignment \cite{wallace2024diffusiondpo}, we replace log-probabilities with a denoising-based implicit reward score for each side of the preference pair. For $a\in\{+,-\}$, where $+$ denotes the preferred sample and $-$ denotes the rejected sample, training evaluates a clean target-part latent $z_{p,0}^a$, sampled noise $\epsilon^a$, noisy latent
$z_{p,t}^a=(1-\sigma_t)z_{p,0}^a+\sigma_t\epsilon^a$, and configured denoising target $u^a$; in our PartCrafter runs, $u^a$ is the reverse-velocity target. Positive-only SFT samples do not define a rejected sample and therefore use only the denoising objective. Synthetic-negative controls construct the rejected sample by explicitly corrupting the target-part latent.
We define
{\small
\begin{equation}
s_\theta(a|x,p,t) =
- \left\|
u^a - \hat{u}_\theta(z_{p,t}^a,t,x,p)
\right\|_2^2 ,
\label{eq:denoise_score}
\end{equation}
}
and analogously define $s_{\mathrm{ref}}$ using the frozen pretrained reference policy. We then compute the policy and reference margins
{\small
\begin{equation}
\begin{aligned}
\Delta_\theta
&= s_\theta(+|x,p,t)-s_\theta(-|x,p,t),\\
\Delta_{\mathrm{ref}}
&= s_{\mathrm{ref}}(+|x,p,t)
-s_{\mathrm{ref}}(-|x,p,t),
\end{aligned}
\label{eq:pref_margin}
\end{equation}
}
which gives the implemented preference term
{\small
\begin{equation}
\mathcal{L}_{\mathrm{pref}}
= - \log \sigma\!\left(\beta(\Delta_\theta-\Delta_{\mathrm{ref}})\right).
\label{eq:pref}
\end{equation}
}
% \SK{What is $\epsilon_{\theta}$?}
Here $\hat{u}_\theta$ is the trainable denoiser's prediction in the configured PartCrafter training-objective space. The objective increases the preferred-over-rejected implicit reward margin after subtracting the corresponding reference-policy margin, so updates target preference gaps not already explained by the pretrained model.

% Previous visible wording: Operationally, branch inputs are represented as clean target-part latents and sampled noise trajectories, then passed through the same trainable DiT with shared weights. Synthetic-negative controls explicitly corrupt the rejected target latent; generated-pair settings use the selected stochastic branches for the same $(x,p)$. The frozen reference DiT evaluates the same branches without gradients. Figure~\ref{fig:dit_preference_training} summarizes the update within the PartCrafter-style backbone.
Operationally, the preferred and rejected samples are noised at a sampled timestep and evaluated by the same trainable DiT. The frozen reference DiT evaluates the same samples without gradients, so the pairwise term depends on the preference margin relative to the pretrained model. Positive-only SFT samples use the denoising term without a rejected sample; the supplementary material provides the full training schematic.

We combine the pairwise term with a standard denoising regularizer on the preferred sample:
\begin{equation}
\mathcal{L}
= \lambda_{\mathrm{pref}}\mathcal{L}_{\mathrm{pref}}
+ \lambda_{\mathrm{reg}}
\left\|
u^+ - \hat{u}_\theta(z_{p,t}^+,t,x,p)
\right\|_2^2 .
\label{eq:total_loss}
\end{equation}
The regularizer is applied only to the preferred sample, limiting drift while the pairwise term separates preferred from rejected samples. Baselines that remove the rejected sample are treated as controls in Sec.~\ref{sec:baseline_controls}.

% Figure moved to supplementary material to reduce main-paper figure footprint.
% \begin{figure*}[t]
% \centering
% \includegraphics[width=0.82\textwidth,keepaspectratio]{figures/figure1_pair_selection_chair_20260707/dit_training_schematic/figure1_right_dit_training_schematic_vertical.png}
% \caption{Preference fine-tuning on the PartCrafter-style DiT backbone. The input image provides feature tokens, while preferred and rejected target-part meshes provide noisy part-latent tokens. The trainable DiT and frozen reference evaluate the same pair; the pairwise loss and denoising regularizer update only the DiT denoiser.}
% \label{fig:dit_preference_training}
% \end{figure*}

% \subsection{Reference-Anchored Margins Beyond Positive-Only SFT}
% Direct SFT and pairwise preference fine-tuning use overlapping supervision but impose different optimization constraints. Direct SFT increases agreement with curated positive targets through a denoising loss and should therefore be competitive when those positives are reliable and close to the evaluation distribution. The pairwise objective instead optimizes a preferred-versus-rejected margin for the same input and target part. In Eq.~\ref{eq:pref}, this margin is measured relative to the frozen pretrained reference; the update therefore favors the preferred branch only when it improves upon the reference-normalized rejected branch. This distinction is most relevant when both candidates preserve object-level structure but differ in local part quality. It also motivates the object-scale study. With too few source objects, the margin signal may be narrow or noisy. With broader curated object coverage, the signal should become less idiosyncratic. If additional object-derived pairs are weakly curated, however, increasing the set size need not improve the model. We therefore treat object-set scale as an empirical variable rather than assuming monotonic gains.

% Moved from Method to Experiments on 2026-07-14.
% \subsection{Curated Candidate Pool}
% We use a curated pretrained-PartCrafter candidate set from PartObjaverse-Tiny \cite{yang2024sampart3d}, with 300 same-input target-part preference examples in the primary metric-assisted setting. Objects are selected after filtering very small parts and retaining visible semantic coverage. Each sample pairs preferred and rejected branches from the frozen pretrained model for the same object and target part. Generation-quality claims are based on fixed evaluations, not on construction labels.
%
% We evaluate three effective pair-set sizes from this pool to separate curated coverage from objective design.
%
% Curated labels define supervision but do not establish generation quality. We therefore use Table~\ref{tab:main_results} as the primary quantitative evidence for model behavior.
%
% The curated setup is intentionally small and inspectable, allowing the construction process to be audited before more automated collection.
%
% \subsection{Curated Pair-Set Size Controls}
% \label{sec:pair_set_scaling}
% We compare 56, 363, and 796 same-input preference pairs, derived from progressively larger source-object pools. Reporting effective pair count reflects that each retained semantic-part row contributes one training preference. The comparison tests whether the same protocol remains informative as coverage increases, while probing whether added pairs become redundant, less reliable, or distributionally mismatched.

\section{Experiments}
The experiments test whether the compact preference signal improves held-out part generation and which factors control the effect. Sec.~\ref{sec:experiment_protocol} gives the data, training, and evaluation details; Sec.~\ref{sec:overall_results} compares the preference model with pretrained PartCrafter and direct SFT; Sec.~\ref{sec:ablation_reliability} analyzes pair construction and reliability; Sec.~\ref{sec:ablation_strength} checks objective strength.

\subsection{Implementation and Evaluation Protocol}
\label{sec:experiment_protocol}
\noindent\textbf{Data and pair counts.}
We use a curated set of pretrained-PartCrafter samples from PartObjaverse-Tiny \cite{yang2024sampart3d}. The primary metric-assisted setting contains 300 same-input target-part preference examples. Objects are selected after filtering very small parts and retaining visible semantic coverage. Each retained example contributes one same-input preference pair for a target part, with preferred and rejected samples generated by the frozen pretrained model. For pair-count controls, we additionally evaluate smaller and larger metric-assisted settings from the same curation pipeline; the human-inspected controls are separate 41- and 7-pair pilots.

\noindent\textbf{Training details.}
All preference models start from the released PartCrafter weights and fine-tune only the denoising transformer; the image encoder, VAE decoder, and frozen reference model remain fixed. Unless otherwise stated, preference runs use AdamW with learning rate $1\times 10^{-4}$, constant warmup for 10 steps, fp16 mixed precision, gradient checkpointing, EMA, batch size four per GPU, $\beta=0.10$, preference weight $\lambda_{\mathrm{pref}}=0.40$, selected-part denoising-regularizer weight $0.20$, and all-part preservation-regularizer weight $0.05$. The direct SFT baseline uses the same backbone and optimizer family but removes the rejected sample and trains on curated positive targets with the original denoising objective.

\noindent\textbf{Evaluation protocol.}
The primary evaluation uses fixed multi-part inputs because single-part cases can make object- and part-level metrics nearly equivalent. We stratify by semantic part count and compare the released pretrained PartCrafter baseline with preference-tuned models and direct SFT.

We run fixed-list evaluations on Objaverse \cite{deitke2023objaverse}, Amazon Berkeley Objects (ABO) \cite{collins2022abo}, and ShapeNet \cite{chang2015shapenet}. The Objaverse split excludes PartObjaverse-Tiny source objects but is the closest evaluation distribution because curation uses PartObjaverse-Tiny. ABO and ShapeNet differ in provenance, categories, and part statistics, so they are treated as transfer checks.
% Previous stronger framing: Single-part cases are reported separately because they are not the primary weak-part setting.
Single-part cases are reported separately because they are less diagnostic of part-level supervision. Tables report CD and FS at object and target-part levels; FS uses a threshold of $0.1$ times the normalized ground-truth bounding-box diagonal. Empty or degenerate target-part outputs are excluded from part-level aggregation under the same protocol. Bold values improve over the corresponding pretrained baseline, and underlines mark the best value within each comparison group.

\subsection{Overall Results and Baseline Controls}
\label{sec:overall_results}
% Full Table 1 moved to supplementary material on 2026-07-14.
\iffalse
\subsection{Main Results: Pair Sources and Reliability Controls}
\begin{table*}[t]
\centering
% Previous caption wording: The Setting column separates the released baseline, correct real preference pairs, noisy pair-order controls, synthetic negative controls, and a current human-annotated pilot. Correct real pairs are metric/ground-truth-assisted preferred/rejected pairs curated from pretrained PartCrafter candidate branches under the same input and target part.
\caption{Main comparison of pair sources and reliability controls. Metric-assisted settings use same-input target-part pairs selected by part-level scores and curation checks. \emph{+crop} removes part-surface regions, \emph{+flip} reverses labels, and \emph{+noise} injects uncurated seed-pair samples.}
\label{tab:main_results}
\resulttablesize
\resulttablespacing{0.5pt}
\begin{tabular*}{\textwidth}{@{\extracolsep{\fill}}llrrrrrrrrrrrr@{}}
\toprule
Method & Setting & \multicolumn{4}{c}{\textbf{Objaverse}} & \multicolumn{4}{c}{\textbf{ABO}} & \multicolumn{4}{c}{\textbf{ShapeNet}} \\
\cmidrule(lr){3-6}\cmidrule(lr){7-10}\cmidrule(lr){11-14}
 & & \multicolumn{2}{c}{Object} & \multicolumn{2}{c}{Part} & \multicolumn{2}{c}{Object} & \multicolumn{2}{c}{Part} & \multicolumn{2}{c}{Object} & \multicolumn{2}{c}{Part} \\
\cmidrule(lr){3-4}\cmidrule(lr){5-6}\cmidrule(lr){7-8}\cmidrule(lr){9-10}\cmidrule(lr){11-12}\cmidrule(lr){13-14}
 & & CD $\downarrow$ & FS $\uparrow$ & CD $\downarrow$ & FS $\uparrow$ & CD $\downarrow$ & FS $\uparrow$ & CD $\downarrow$ & FS $\uparrow$ & CD $\downarrow$ & FS $\uparrow$ & CD $\downarrow$ & FS $\uparrow$ \\
\midrule
Baseline & Pretrained & 0.2724 & 0.5680 & 0.2875 & 0.5453 & 0.1520 & 0.7806 & 0.3204 & 0.4888 & 0.1814 & 0.7130 & 0.3262 & 0.4457 \\
\midrule
% Previous row label: GT-guided auto.
\multirow{3}{*}{Metric-assisted} & 56 pairs & \metriccd{0.2626}{0.2724}{0.1986} & \metricf{0.5328}{0.5680}{0.6958} & \metriccd{0.3157}{0.2875}{0.2542} & \metricf{0.4545}{0.5453}{0.5725} & \metriccd{0.3000}{0.1520}{0.2233} & \metricf{0.4528}{0.7806}{0.6560} & \metriccd{0.3722}{0.3204}{0.3104} & \metricf{0.3437}{0.4888}{0.4707} & \metriccd{0.2843}{0.1814}{0.1833} & \metricf{0.4442}{0.7130}{0.7118} & \metriccd{0.3612}{0.3262}{0.2872} & \metricf{0.3500}{0.4457}{0.4688} \\
 & 363 pairs & \metriccd{0.1986}{0.2724}{0.1986} & \metricf{0.6958}{0.5680}{0.6958} & \metriccd{0.2542}{0.2875}{0.2542} & \metricf{0.5725}{0.5453}{0.5725} & \metriccd{0.2233}{0.1520}{0.2233} & \metricf{0.6560}{0.7806}{0.6560} & \metriccd{0.3104}{0.3204}{0.3104} & \metricf{0.4707}{0.4888}{0.4707} & \metriccd{0.2099}{0.1814}{0.1833} & \metricf{0.6500}{0.7130}{0.7118} & \metriccd{0.2872}{0.3262}{0.2872} & \metricf{0.4688}{0.4457}{0.4688} \\
 & \textbf{796 pairs} & \metriccd{0.2625}{0.2724}{0.1986} & \metricf{0.5881}{0.5680}{0.6958} & \metriccd{0.2716}{0.2875}{0.2542} & \metricf{0.5707}{0.5453}{0.5725} & \metriccd{0.1382}{0.1520}{0.1382} & \metricf{0.8072}{0.7806}{0.8072} & \metriccd{0.3215}{0.3204}{0.3104} & \metricf{0.4729}{0.4888}{0.4729} & \metriccd{0.1833}{0.1814}{0.1833} & \metricf{0.7118}{0.7130}{0.7118} & \metriccd{0.3344}{0.3262}{0.2872} & \metricf{0.4396}{0.4457}{0.4688} \\
\midrule
% Previous row label: Human-annotated.
\multirow{2}{*}{Human-inspected} & 41 pairs & \metriccd{0.5127}{0.2724}{0.2603} & \metricf{0.1745}{0.5680}{0.5854} & \metriccd{0.5616}{0.2875}{0.2805} & \metricf{0.1285}{0.5453}{0.5573} & \metriccd{0.4440}{0.1520}{0.1400} & \metricf{0.2215}{0.7806}{0.8021} & \metriccd{0.5293}{0.3204}{0.3096} & \metricf{0.1531}{0.4888}{0.5006} & \metriccd{0.4942}{0.1814}{0.1725} & \metricf{0.1579}{0.7130}{0.7386} & \metriccd{0.5557}{0.3262}{0.3267} & \metricf{0.1284}{0.4457}{0.4493} \\
 & 7 pairs & \metriccd{0.2603}{0.2724}{0.2603} & \metricf{0.5854}{0.5680}{0.5854} & \metriccd{0.2805}{0.2875}{0.2805} & \metricf{0.5573}{0.5453}{0.5573} & \metriccd{0.1400}{0.1520}{0.1400} & \metricf{0.8021}{0.7806}{0.8021} & \metriccd{0.3096}{0.3204}{0.3096} & \metricf{0.5006}{0.4888}{0.5006} & \metriccd{0.1725}{0.1814}{0.1725} & \metricf{0.7386}{0.7130}{0.7386} & \metriccd{0.3267}{0.3262}{0.3267} & \metricf{0.4493}{0.4457}{0.4493} \\
\midrule
\multirow{2}{*}{Synthetic neg.} & +crop 30\% & \metriccd{0.3529}{0.2724}{0.3529} & \metricf{0.4127}{0.5680}{0.4127} & \metriccd{0.3821}{0.2875}{0.3821} & \metricf{0.3398}{0.5453}{0.3398} & \metriccd{0.3319}{0.1520}{0.3319} & \metricf{0.4355}{0.7806}{0.4355} & \metriccd{0.4145}{0.3204}{0.4145} & \metricf{0.3072}{0.4888}{0.3072} & \metriccd{0.3681}{0.1814}{0.3514} & \metricf{0.3531}{0.7130}{0.3717} & \metriccd{0.4490}{0.3262}{0.4346} & \metricf{0.2541}{0.4457}{0.2631} \\
 & +crop 50\% & \metriccd{0.3636}{0.2724}{0.3529} & \metricf{0.3721}{0.5680}{0.4127} & \metriccd{0.4222}{0.2875}{0.3821} & \metricf{0.2836}{0.5453}{0.3398} & \metriccd{0.3610}{0.1520}{0.3319} & \metricf{0.3950}{0.7806}{0.4355} & \metriccd{0.4280}{0.3204}{0.4145} & \metricf{0.2945}{0.4888}{0.3072} & \metriccd{0.3514}{0.1814}{0.3514} & \metricf{0.3717}{0.7130}{0.3717} & \metriccd{0.4346}{0.3262}{0.4346} & \metricf{0.2631}{0.4457}{0.2631} \\
\midrule
\multirow{3}{*}{Pair-order flip} & +flip 10\% & \metriccd{0.3848}{0.2724}{0.3628} & \metricf{0.3238}{0.5680}{0.3709} & \metriccd{0.4336}{0.2875}{0.4267} & \metricf{0.2690}{0.5453}{0.2858} & \metriccd{0.3678}{0.1520}{0.3678} & \metricf{0.3957}{0.7806}{0.3957} & \metriccd{0.4562}{0.3204}{0.4562} & \metricf{0.2672}{0.4888}{0.2672} & \metriccd{0.3722}{0.1814}{0.3722} & \metricf{0.3478}{0.7130}{0.3478} & \metriccd{0.4289}{0.3262}{0.4289} & \metricf{0.2645}{0.4457}{0.2645} \\
 & +flip 30\% & \metriccd{0.3628}{0.2724}{0.3628} & \metricf{0.3709}{0.5680}{0.3709} & \metriccd{0.4267}{0.2875}{0.4267} & \metricf{0.2858}{0.5453}{0.2858} & \metriccd{0.4026}{0.1520}{0.3678} & \metricf{0.3161}{0.7806}{0.3957} & \metriccd{0.4783}{0.3204}{0.4562} & \metricf{0.2384}{0.4888}{0.2672} & \metriccd{0.3804}{0.1814}{0.3722} & \metricf{0.3102}{0.7130}{0.3478} & \metriccd{0.4560}{0.3262}{0.4289} & \metricf{0.2434}{0.4457}{0.2645} \\
 & +flip 50\% & \metriccd{0.3738}{0.2724}{0.3628} & \metricf{0.3487}{0.5680}{0.3709} & \metriccd{0.4480}{0.2875}{0.4267} & \metricf{0.2579}{0.5453}{0.2858} & \metriccd{0.3996}{0.1520}{0.3678} & \metricf{0.3247}{0.7806}{0.3957} & \metriccd{0.4822}{0.3204}{0.4562} & \metricf{0.2192}{0.4888}{0.2672} & \metriccd{0.3833}{0.1814}{0.3722} & \metricf{0.3178}{0.7130}{0.3478} & \metriccd{0.4650}{0.3262}{0.4289} & \metricf{0.2194}{0.4457}{0.2645} \\
% & +flip 70\% & 0.3470 & 0.3862 & 0.4100 & 0.2995 & 0.3016 & 0.4747 & 0.4049 & 0.3144 & 0.3790 & 0.3189 & 0.4542 & 0.2387 \\
\midrule
\multirow{3}{*}{Noisy samples} & +noise 10\% & \metriccd{0.2461}{0.2724}{0.2461} & \metricf{0.5330}{0.5680}{0.5330} & \metriccd{0.3122}{0.2875}{0.3122} & \metricf{0.4319}{0.5453}{0.4319} & \metriccd{0.2267}{0.1520}{0.2267} & \metricf{0.6041}{0.7806}{0.6041} & \metriccd{0.3287}{0.3204}{0.3287} & \metricf{0.4440}{0.4888}{0.4440} & \metriccd{0.2432}{0.1814}{0.2432} & \metricf{0.5240}{0.7130}{0.5240} & \metriccd{0.3227}{0.3262}{0.3227} & \metricf{0.3975}{0.4457}{0.3975} \\
 & +noise 30\% & \metriccd{0.3542}{0.2724}{0.2461} & \metricf{0.3554}{0.5680}{0.5330} & \metriccd{0.4398}{0.2875}{0.3122} & \metricf{0.2364}{0.5453}{0.4319} & \metriccd{0.3549}{0.1520}{0.2267} & \metricf{0.2719}{0.7806}{0.6041} & \metriccd{0.5075}{0.3204}{0.3287} & \metricf{0.1812}{0.4888}{0.4440} & \metriccd{0.3982}{0.1814}{0.2432} & \metricf{0.2512}{0.7130}{0.5240} & \metriccd{0.5219}{0.3262}{0.3227} & \metricf{0.1458}{0.4457}{0.3975} \\
 & +noise 50\% & \metriccd{0.3635}{0.2724}{0.2461} & \metricf{0.3797}{0.5680}{0.5330} & \metriccd{0.4009}{0.2875}{0.3122} & \metricf{0.3043}{0.5453}{0.4319} & \metriccd{0.3147}{0.1520}{0.2267} & \metricf{0.3820}{0.7806}{0.6041} & \metriccd{0.4057}{0.3204}{0.3287} & \metricf{0.2755}{0.4888}{0.4440} & \metriccd{0.4225}{0.1814}{0.2432} & \metricf{0.2552}{0.7130}{0.5240} & \metriccd{0.4958}{0.3262}{0.3227} & \metricf{0.1892}{0.4457}{0.3975} \\
\bottomrule
\end{tabular*}
\end{table*}
\fi

% Previous wording: The correct real-pair rows compare metric/ground-truth-assisted preferred/rejected pairs curated from pretrained PartCrafter candidate branches at three effective pair-set sizes.
% Previous wording: Table~\ref{tab:main_results} groups the main comparison by pair source and reliability control. The metric-assisted rows show that generated same-input target-part pairs remain useful as pair coverage increases, but the trend is not monotonic: the $\sim$350-pair setting is strongest for several part-level metrics, while the completed $\sim$800-pair checkpoint improves all Objaverse metrics and ABO object-level metrics over the pretrained baseline. This supports the data-quality view that additional pairs help only when their labels remain reliable and distributionally aligned. The human-inspected row is a small pilot from 41 effective same-input pairs after filtering invalid selections, so we treat it as an annotation-source control rather than as a final human-labeling result. Synthetic negatives provide a usable but weaker augmentation signal than generated-pair curation. Pair-order flips and noisy-sample injection further show that reliable preferred--rejected ordering is essential: +noise 10\% is the least damaging noisy-sample setting, but 30\% and 50\% injection sharply degrade Part FS across all three datasets.
% Previous visible wording: Table~\ref{tab:main_results} groups results by pair source and reliability control. Metric-assisted preferences require reliable coverage: the $\sim$50-pair setting improves Objaverse object CD but not part metrics, while $\sim$350 pairs gives the strongest matched-distribution Objaverse part results. The completed $\sim$800-pair checkpoint remains above the pretrained baseline on all Objaverse metrics and improves ABO object-level CD and FS, but cross-dataset part gains are weaker. This supports a data-quality interpretation rather than a scale-only view: more pairs help only when added labels remain reliable and aligned.

% Previous visible wording: The controls show which additional supervision is less effective. The human-inspected row is a small 41-pair annotation-source pilot. Synthetic crop negatives are weaker than generated-pair curation, suggesting that handcrafted missing-part negatives do not replace real same-input alternatives. Pair-order flips degrade results, and noisy-sample injection is unstable: 10\% is least damaging, but noisy settings reduce FS relative to the pretrained baseline and clean metric-assisted rows. Overall, curated preferences help most in the closest evaluation setting, with label reliability and distribution match central to the benefit.


\label{sec:baseline_controls}
\begin{table*}[t]
\centering
\caption{Overall results and part-count controls. Bold ``Ours'' denotes the reported DPO-style preference model. SFT uses curated preferred samples without a rejected sample. The all-parts average combines underlying fixed-list samples across buckets.}
\label{tab:baseline_controls}
\resulttablesize
\resulttablespacing{0.5pt}
\begin{tabular*}{\textwidth}{@{\extracolsep{\fill}}llrrrrrrrrrrrr@{}}
\toprule
Part Numbers & Method & \multicolumn{4}{c}{\textbf{Objaverse}} & \multicolumn{4}{c}{\textbf{ABO}} & \multicolumn{4}{c}{\textbf{ShapeNet}} \\
\cmidrule(lr){3-6}\cmidrule(lr){7-10}\cmidrule(lr){11-14}
 & & \multicolumn{2}{c}{Object} & \multicolumn{2}{c}{Part} & \multicolumn{2}{c}{Object} & \multicolumn{2}{c}{Part} & \multicolumn{2}{c}{Object} & \multicolumn{2}{c}{Part} \\
\cmidrule(lr){3-4}\cmidrule(lr){5-6}\cmidrule(lr){7-8}\cmidrule(lr){9-10}\cmidrule(lr){11-12}\cmidrule(lr){13-14}
 & & CD $\downarrow$ & FS $\uparrow$ & CD $\downarrow$ & FS $\uparrow$ & CD $\downarrow$ & FS $\uparrow$ & CD $\downarrow$ & FS $\uparrow$ & CD $\downarrow$ & FS $\uparrow$ & CD $\downarrow$ & FS $\uparrow$ \\
\midrule
\multirow{3}{*}{1 part} & Pretrained & \metricbest{0.2382}{0.2184} & \metricbest{0.6366}{0.6603} & \metricbest{0.2096}{0.1890} & \metricbest{0.6949}{0.7273} & \metricbest{0.0978}{0.0978} & \metricbest{0.8998}{0.9066} & \metricbest{0.0975}{0.0975} & \metricbest{0.9007}{0.9007} & \metricbest{0.2092}{0.2092} & \metricbest{0.6626}{0.6626} & \metricbest{0.2094}{0.2094} & \metricbest{0.6632}{0.6632} \\
 & SFT & \metriccd{0.2243}{0.2382}{0.2184} & \metricf{0.6441}{0.6366}{0.6603} & \metriccd{0.1890}{0.2096}{0.1890} & \metricf{0.7273}{0.6949}{0.7273} & \metriccd{0.0991}{0.0978}{0.0978} & \metricf{0.9066}{0.8998}{0.9066} & \metriccd{0.1043}{0.0975}{0.0975} & \metricf{0.8935}{0.9007}{0.9007} & \metriccd{0.2162}{0.2092}{0.2092} & \metricf{0.6411}{0.6626}{0.6626} & \metriccd{0.2161}{0.2094}{0.2094} & \metricf{0.6401}{0.6632}{0.6632} \\
 & \textbf{Ours} & \metriccd{0.2184}{0.2382}{0.2184} & \metricf{0.6603}{0.6366}{0.6603} & \metriccd{0.2044}{0.2096}{0.1890} & \metricf{0.7020}{0.6949}{0.7273} & \metriccd{0.1007}{0.0978}{0.0978} & \metricf{0.8883}{0.8998}{0.9066} & \metriccd{0.1015}{0.0975}{0.0975} & \metricf{0.8898}{0.9007}{0.9007} & \metriccd{0.2276}{0.2092}{0.2092} & \metricf{0.6147}{0.6626}{0.6626} & \metriccd{0.2294}{0.2094}{0.2094} & \metricf{0.6000}{0.6632}{0.6632} \\
\midrule
\multirow{3}{*}{2--4 parts} & Pretrained & \metricbest{0.2631}{0.2491} & \metricbest{0.5909}{0.6121} & \metricbest{0.2574}{0.2485} & \metricbest{0.5986}{0.6237} & \metricbest{0.1520}{0.1382} & \metricbest{0.7806}{0.8087} & \metricbest{0.3204}{0.3062} & \metricbest{0.4888}{0.4889} & \metricbest{0.1828}{0.1807} & \metricbest{0.7107}{0.7179} & \metricbest{0.3115}{0.3099} & \metricbest{0.4689}{0.4748} \\
 & SFT & \metriccd{0.2568}{0.2631}{0.2491} & \metricf{0.5945}{0.5909}{0.6121} & \metriccd{0.2620}{0.2574}{0.2485} & \metricf{0.5955}{0.5986}{0.6237} & \metriccd{0.1409}{0.1520}{0.1382} & \metricf{0.8087}{0.7806}{0.8087} & \metriccd{0.3062}{0.3204}{0.3062} & \metricf{0.4889}{0.4888}{0.4889} & \metriccd{0.1807}{0.1828}{0.1807} & \metricf{0.7178}{0.7107}{0.7179} & \metriccd{0.3099}{0.3115}{0.3099} & \metricf{0.4748}{0.4689}{0.4748} \\
 & \textbf{Ours} & \metriccd{0.2491}{0.2631}{0.2491} & \metricf{0.6121}{0.5909}{0.6121} & \metriccd{0.2485}{0.2574}{0.2485} & \metricf{0.6237}{0.5986}{0.6237} & \metriccd{0.1382}{0.1520}{0.1382} & \metricf{0.8072}{0.7806}{0.8087} & \metriccd{0.3215}{0.3204}{0.3062} & \metricf{0.4729}{0.4888}{0.4889} & \metriccd{0.1816}{0.1828}{0.1807} & \metricf{0.7179}{0.7107}{0.7179} & \metriccd{0.3199}{0.3115}{0.3099} & \metricf{0.4712}{0.4689}{0.4748} \\
\midrule
\multirow{3}{*}{5--8 parts} & Pretrained & \metricbest{0.2865}{0.2813} & \metricbest{0.5332}{0.5520} & \metricbest{0.3333}{0.3062} & \metricbest{0.4646}{0.4912} & -- & -- & -- & -- & \metricbest{0.1795}{0.1677} & \metricbest{0.7160}{0.7391} & \metricbest{0.3456}{0.3456} & \metricbest{0.4149}{0.4149} \\
 & SFT & \metriccd{0.2813}{0.2865}{0.2813} & \metricf{0.5463}{0.5332}{0.5520} & \metriccd{0.3197}{0.3333}{0.3062} & \metricf{0.4846}{0.4646}{0.4912} & -- & -- & -- & -- & \metriccd{0.1677}{0.1795}{0.1677} & \metricf{0.7391}{0.7160}{0.7391} & \metriccd{0.3497}{0.3456}{0.3456} & \metricf{0.4006}{0.4149}{0.4149} \\
 & \textbf{Ours} & \metriccd{0.2824}{0.2865}{0.2813} & \metricf{0.5520}{0.5332}{0.5520} & \metriccd{0.3062}{0.3333}{0.3062} & \metricf{0.4912}{0.4646}{0.4912} & -- & -- & -- & -- & \metriccd{0.1854}{0.1795}{0.1677} & \metricf{0.7037}{0.7160}{0.7391} & \metriccd{0.3537}{0.3456}{0.3456} & \metricf{0.3975}{0.4149}{0.4149} \\

\midrule
\multirow{3}{*}{\textbf{ avg.}} & Pretrained & \metricbest{0.2724}{0.2625} & \metricbest{0.5680}{0.5881} & \metricbest{0.2875}{0.2716} & \metricbest{0.5453}{0.5707} & \metricbest{0.1520}{0.1382} & \metricbest{0.7806}{0.8087} & \metricbest{0.3204}{0.3062} & \metricbest{0.4888}{0.4889} & \metricbest{0.1814}{0.1751} & \metricbest{0.7130}{0.7269} & \metricbest{0.3262}{0.3262} & \metricbest{0.4457}{0.4457} \\
 & SFT & \metriccd{0.2667}{0.2724}{0.2625} & \metricf{0.5751}{0.5680}{0.5881} & \metriccd{0.2852}{0.2875}{0.2716} & \metricf{0.5509}{0.5453}{0.5707} & \metriccd{0.1409}{0.1520}{0.1382} & \metricf{0.8087}{0.7806}{0.8087} & \metriccd{0.3062}{0.3204}{0.3062} & \metricf{0.4889}{0.4888}{0.4889} & \metriccd{0.1751}{0.1814}{0.1751} & \metricf{0.7269}{0.7130}{0.7269} & \metriccd{0.3270}{0.3262}{0.3262} & \metricf{0.4430}{0.4457}{0.4457} \\
 & \textbf{Ours} & \metriccd{0.2625}{0.2724}{0.2625} & \metricf{0.5881}{0.5680}{0.5881} & \metriccd{0.2716}{0.2875}{0.2716} & \metricf{0.5707}{0.5453}{0.5707} & \metriccd{0.1382}{0.1520}{0.1382} & \metricf{0.8072}{0.7806}{0.8087} & \metriccd{0.3215}{0.3204}{0.3062} & \metricf{0.4729}{0.4888}{0.4889} & \metriccd{0.1833}{0.1814}{0.1751} & \metricf{0.7118}{0.7130}{0.7269} & \metriccd{0.3344}{0.3262}{0.3262} & \metricf{0.4396}{0.4457}{0.4457} \\
\bottomrule
\end{tabular*}
\end{table*}

Table~\ref{tab:baseline_controls} stratifies fixed-list evaluation by part count across Objaverse, ABO, and ShapeNet. Higher part counts require more local components to be generated and separated.
% Previous stronger framing: The single-part rows are auxiliary sanity controls and are not used as primary evidence for weak-part decomposition.
The single-part entries are auxiliary sanity controls and are less diagnostic of the part-level supervision question. The all-parts average is computed from the underlying fixed-split samples, not by averaging displayed bucket means.
% Previous stronger framing: The primary weak-part claim is evaluated through same-dataset, same-bucket comparisons against the pretrained PartCrafter baseline and the direct SFT control.
Primary quantitative comparisons are made within the same dataset and part-count bucket against the pretrained baseline and direct SFT control.

% Previous wording: The strongest improvements appear on Objaverse, which is also the evaluation distribution closest to the PartObjaverse-Tiny source used for pair curation. This pattern is consistent with a near-distribution adaptation effect: the preference pairs are curated from source objects whose local part statistics are closer to Objaverse-style objects. The weaker or less consistent gains on ABO and ShapeNet should therefore be read as evidence about distribution shift and cross-dataset transfer, not as the same operating condition as Objaverse.
The strongest improvements appear on Objaverse, the evaluation distribution closest to PartObjaverse-Tiny. Weaker or less consistent ABO and ShapeNet gains should therefore be read as evidence about distribution shift and transfer, not as failures under the same operating condition.

\subsection{Ablation: Pair Construction and Reliability}
\label{sec:ablation_reliability}

% Previous wording: The \emph{+flip} rows show that noisy pair order weakens part-level metrics relative to the clean correct real-pair setting.
% Previous visible wording: The ablations in Table~\ref{tab:main_results} separate curated pair quality from simply adding rows. Pair-order flips, noisy-sample injection, and synthetic cropping underperform metric-assisted generated pairs, suggesting that the useful signal comes from reliable local preferences rather than generic exposure to extra samples.
The full pair-source and reliability table is moved to the supplementary material. Its ablations separate curated pair quality from simply adding more examples: pair-order flips, noisy-sample injection, and synthetic cropping underperform metric-assisted generated pairs, suggesting that the useful signal comes from reliable local preferences rather than generic exposure to extra samples. The seven-pair screened pilot improves all four Objaverse and ABO metrics and three of four ShapeNet metrics; ShapeNet Part CD is effectively unchanged at 0.3267 versus 0.3262. The contrast with the weaker 41-pair pilot is consistent with pair selection and reliability mattering, but it does not isolate pair count because the pilots differ in source data, screening, and training settings.

\subsection{Ablation: Preference Strength}
\label{sec:ablation_strength}

% Full qualitative figure moved to supplementary material; a compact Ours-only crop remains below.
% \begin{figure}[t]
% \centering
% \includegraphics[width=0.98\linewidth,keepaspectratio]{figures/figure4_relayout_assets_20260709/figure4_selected_M01_M05_X03_layout_two_row_zoom.png}
% \caption{Qualitative multi-view comparisons. Each row uses the same input image and part count for the pretrained and preference-tuned models. Matched views show object consistency and part separation; zoom-in crops use the same highlighted local region for both methods.}
% \label{fig:prediction_examples}
% \end{figure}

\begin{table}[t]
\centering
\caption{Preference-strength ablation on the ShapeNet multi-part fixed list. Values are averaged over the fixed list; object- and part-level optima differ.}
\label{tab:beta_ablation}
\resulttablespacing{3pt}
{\resulttablesize
\begin{tabular*}{0.98\columnwidth}{@{\extracolsep{\fill}}lrrrr@{}}
\toprule
Setting & \multicolumn{4}{c}{\textbf{Fixed-list average}} \\
\cmidrule(lr){2-5}
 & \multicolumn{2}{c}{Object} & \multicolumn{2}{c}{Part} \\
\cmidrule(lr){2-3}\cmidrule(lr){4-5}
$\beta$ & CD $\downarrow$ & FS $\uparrow$ & CD $\downarrow$ & FS $\uparrow$ \\
\midrule
0.05 & \partialmetric{0.3003} & \partialmetric{0.4293} & \partialmetric{0.3872} & \partialmetric{0.2931} \\
0.10 & \bestpartialmetric{0.1776} & \bestpartialmetric{0.7237} & \partialmetric{0.3288} & \partialmetric{0.4513} \\
\textbf{0.20} & \partialmetric{0.2249} & \partialmetric{0.6194} & \bestpartialmetric{0.3049} & \bestpartialmetric{0.4564} \\
0.40 & \partialmetric{0.2755} & \partialmetric{0.5033} & \partialmetric{0.3580} & \partialmetric{0.3368} \\
\bottomrule
\end{tabular*}
}
\end{table}
\vspace{-0.5em}

Table~\ref{tab:beta_ablation} checks preference-strength sensitivity. On ShapeNet multi-part inputs, $\beta=0.10$ gives the strongest object-level CD and FS, while $\beta=0.20$ gives the best part-level CD and FS. The larger $\beta=0.40$ degrades both levels, suggesting that an overly strong pairwise term can destabilize generation rather than simply strengthening preference adaptation.

A controlled ShapeNet fixed-list evaluation further examines round-to-round behavior. Against the released pretrained baseline, the reported metric-assisted setting improves Part FS from \partialmetric{0.37} to \partialmetric{0.40} after one round and \partialmetric{0.42} after two rounds; Part CD decreases from \partialmetric{0.39} to \partialmetric{0.36}. Cross-dataset results in Table~\ref{tab:baseline_controls} and the supplementary pair-source table remain the main evidence for transfer behavior.

Direct SFT is retained as a strong baseline because curated positives can provide a strong adaptation signal, especially near the target distribution. We therefore do not claim that pairwise preference learning universally dominates SFT or every transfer setting.
% Previous stronger framing: curated part-level preferences can improve weak semantic parts over the pretrained generator under controlled multi-part settings.
The supported claim is narrower: curated part-level preferences provide a measurable local signal under controlled multi-part settings, with benefit depending on pair coverage, pair quality, objective balance, and dataset shift.

\subsection{Qualitative Illustration}
Figure~\ref{fig:prediction_examples_compact} keeps one Ours-only qualitative example from the full comparison sheet to show the visual form of the local improvements. The complete multi-view PartCrafter/Ours comparison is included in the supplementary material.

\begin{figure}[t]
\centering
\includegraphics[width=0.16\linewidth,trim=1590 1385 1231 575,clip]{figures/figure4_relayout_assets_20260709/figure4_selected_M01_M05_X03_layout_two_row_zoom.png}\hfill
\includegraphics[width=0.80\linewidth,trim=1980 1250 31 735,clip]{figures/figure4_relayout_assets_20260709/figure4_selected_M01_M05_X03_layout_two_row_zoom.png}
\caption{Compact qualitative example. Left: input image. Right: Ours-only views and zoom-in crop from the full qualitative comparison.}
\label{fig:prediction_examples_compact}
\end{figure}

\FloatBarrier

\section{Discussion and Limitations}
This work explores pairwise preferences for part-level supervision. Curated pairs make the signal inspectable but are not yet a scalable collection strategy. Direct SFT remains competitive when curated preferred samples match the evaluation distribution, so preference learning is complementary rather than a replacement for positive-target fine-tuning. The evidence is most relevant to multi-part object inputs and should be read as a controlled fixed-split study, not a statistically powered multi-seed benchmark. Scene-style inputs are outside scope because pairs are selected from object-centric PartObjaverse-Tiny samples and do not encode room layout, multi-object interaction, or scene-level correspondence. Other limitations include pair scalability, invalid target-part outputs, object-quality preservation, scene-aware preferences, and correspondence-aware feedback. Future pair mining could calibrate additional diagnostics, such as boundary quality, part overlap, and multiview consistency, after validation on larger multi-part subsets.

\FloatBarrier

\section{Conclusion}
% Previous stronger framing: We presented a controlled study of curated part-level preferences for weak semantic parts in part-aware 3D generation.
We presented a controlled study of curated part-level preferences as lightweight supervision for part-aware 3D generation. The effect is measurable but distribution-dependent. On the closest held-out Objaverse evaluation, the strongest metric-assisted setting improves over the pretrained baseline: Part CD decreases from 0.2875 to 0.2542 and Part FS increases from 0.5453 to 0.5725.
% Previous stronger framing: These gains indicate that same-input pairwise preferences can provide a useful lightweight signal for correcting local part failures without changing the base generator architecture.
These gains indicate that same-input pairwise preferences can provide a useful local supervision signal without changing the base generator architecture.

Control results do not show universal dominance over direct SFT. On ABO and ShapeNet, SFT is often stronger for part-level metrics, and higher-part-count ShapeNet cases remain challenging. This is consistent with the data distribution: curated pairs come from PartObjaverse-Tiny, making Objaverse nearest, while ABO and ShapeNet are shifted transfer settings.
% Previous stronger framing: curated pairwise preferences can improve weak semantic parts when the preference pairs align with the target distribution.
Noisy-sample controls sharpen this conclusion: low-rate random seed-pair injection is least harmful, but it does not substitute for curated same-input preferences and becomes damaging at higher rates. Overall, curated preferences can improve local part quality when pair sources are reliable and distributionally aligned, while transfer to shifted datasets remains limited. Future work should replace manual preference construction with scalable pair mining, improve object-level preservation, and incorporate scene- or correspondence-aware feedback.

\FloatBarrier

{\scriptsize
\bibliographystyle{ieee}
\bibliography{references}
}

\end{document}
